\documentclass[12pt]{article}
\usepackage[utf8]{inputenc}
\usepackage[T1]{fontenc}
\usepackage{hyperref}
\title{Chazumba Mixtec Language Resources for K-12 Curriculum Development}
\author{Compiled by Claude Sonnet 5.5 (claude-sonnet-5-5) for Andr\'as Kornai}
\date{}
\begin{document}
\maketitle

\section{Introduction}
Chazumba Mixtec (ISO 639-3 xtb; Glottocode chaz1235; also called Da'an davi or Mixteco de Chazumba) is an Otomanguean language spoken along the Puebla--Oaxaca border. Sources give about 3,800 speakers from the year 2000 (Wikipedia, Endangered Languages Project), while Joshua Project lists 5,100. The searches found no dictionary, grammar, primer or course for the language itself.

\section{Primary Educational Resources}
No primers, textbooks or school programmes for Chazumba Mixtec were found.

\section{Linguistic and Reference Materials}
\begin{itemize}
\item Wikipedia overview of Chazumba Mixtec, with location and speaker figures: \url{https://en.wikipedia.org/wiki/Chazumba_Mixtec}
\item Glottolog entry, which places it in the Northern Baja Mixtec group: \url{https://glottolog.org/resource/languoid/id/chaz1235.bigmap.html}
\item Endangered Languages Project entry: \url{https://endangeredlanguages.com/node/26669}
\item Arqueolog\'ia Mexicana PDF on Puebla Mixtec, which includes a dictionary (Valinas) and discusses dialect groupings and endangerment. It concerns the Puebla Mixtec varieties in general and is not a Chazumba dictionary: \url{https://arqueologiamexicana.mx/sites/default/files/banco_imagenes/mixteco_de_puebla_diccionario_valinas.pdf}
\end{itemize}

\section{Religious and Community Resources}
Joshua Project lists the scripture status as unspecified, the written status as unknown, and no Audio Bible or Jesus Film: \url{https://joshuaproject.net/index.php/languages/xtb}

\section{Digital and Audio Resources}
Global Recordings Network reports no audio recordings in this language: \url{https://www.globalrecordings.net/en/language/xtb}

\section{Curriculum Development Context}
Practically no teaching material exists in the search results, and the language is used mainly by older adults. Any curriculum would depend on new documentation with elder speakers.
\end{document}